\documentclass{article}
\usepackage{ctex}
\begin{document}
\ExplSyntaxOn
% 定义条件句，函数名不直接带 p,T,F 字样
% 而是由一个 n 参数指定（逗号分隔）
% 这里的 p 是 “_p” 断言， 和 cs_new:Npn 里的 p 参数意义不一样
\prg_new_conditional:Nnn \my_if_ge_zero:n { p , T , F, TF }
  { 
    % 函数体内只需要返回逻辑真与假即可
    % 其余由 \prg_new_conditional:Nnn 函数处理
    \int_compare:nNnTF
      { #1 } > { 0 }
      { \prg_return_true: }
      { \prg_return_false: }
  }
% 此时定义了如下四个函数变体
\my_if_ge_zero:nTF { 1 } { Y } { N } \par
\my_if_ge_zero:nT { 2 } { 大于零 } \par
\my_if_ge_zero:nF { 3 } { 小于零 } \par
\my_if_ge_zero_p:n { 4 } \par
\ExplSyntaxOff
\end{document}